% ch03_extensions.tex — Metric Extensions
\chapter{Metric Extensions}\label{ch:extensions}

\section{Introduction}

Banach's theorem is elegant, but its hypothesis---a strict, uniform contraction---is sometimes too restrictive. Starting in the 1960s, a whole community of mathematicians embarked on a quest: weaken the assumptions while preserving the conclusion. Edelstein showed that a \emph{local} contraction suffices if the space is compact. Kannan (1968) discovered that a fixed point can exist even when $T$ is not continuous---a surprise that profoundly changed our understanding of this phenomenon. Chatterjea proposed his own variant, \'Ciri\'c unified these results in a general framework, and Caristi introduced an entirely different approach based on decreasing energy functions.

Each extension illuminates a different facet of the mechanism that forces a fixed point to exist, and together they paint a remarkably rich theory.

\section{Edelstein's theorem}

\begin{theorem}[Edelstein, 1962]\label{thm:edelstein}
Let $(X, d)$ be a compact metric space and $T : X \to X$ a weak contraction,
i.e.,
\[
    d(T(x), T(y)) < d(x, y) \quad \forall x, y \in X, \; x \neq y.
\]
Then $T$ has a unique fixed point, and for every $x_0 \in X$, the sequence
$(T^n(x_0))$ converges to this fixed point.
\end{theorem}

\begin{proof}
Define $\varphi : X \to [0, \infty)$ by $\varphi(x) = d(x, Tx)$. This function
is continuous on the compact space $X$, so it attains its minimum at some
$x^* \in X$.

Suppose $\varphi(x^*) > 0$, i.e., $x^* \neq Tx^*$. Then
\[
    \varphi(Tx^*) = d(Tx^*, T^2 x^*) < d(x^*, Tx^*) = \varphi(x^*),
\]
contradicting the minimality of $\varphi(x^*)$. Hence $\varphi(x^*) = 0$ and
$Tx^* = x^*$.

Uniqueness is immediate: if $y^*$ is another fixed point, then
$d(x^*, y^*) = d(Tx^*, Ty^*) < d(x^*, y^*)$, a contradiction.

For convergence, let $x_0 \in X$ and $x_n = T^n(x_0)$. The sequence
$(d(x_n, x^*))$ is decreasing since
$d(x_{n+1}, x^*) = d(Tx_n, Tx^*) < d(x_n, x^*)$ when $x_n \neq x^*$.
By compactness, $(x_n)$ has a convergent subsequence $x_{n_k} \to z$.
Then $d(z, x^*) = \lim d(x_{n_k}, x^*) = \inf_n d(x_n, x^*)$. If $z \neq x^*$,
$d(Tz, x^*) < d(z, x^*)$, but $d(x_{n_k + 1}, x^*) \to d(Tz, x^*)$,
contradicting the fact that $(d(x_n, x^*))$ is decreasing with limit
$d(z, x^*)$. Hence $z = x^*$ and $x_n \to x^*$.
\end{proof}

\begin{warning}[title=\textbf{Compactness is essential}]
Without compactness, Edelstein's theorem fails. The example $T(x) = x + 1/x$ on
$[1, +\infty)$ (a complete but noncompact space) is a weak contraction with no
fixed point (see Chapter~2).
\end{warning}

\section{Kannan contractions}

\begin{definition}[Kannan mapping]
Let $(X, d)$ be a metric space. A mapping $T : X \to X$ is a \emph{Kannan
contraction} if there exists $\alpha \in [0, 1/2)$ such that
\[
    d(Tx, Ty) \leq \alpha \bigl[ d(x, Tx) + d(y, Ty) \bigr] \quad \forall x, y \in X.
\]
\end{definition}

\begin{remark}
Unlike a Banach contraction, a Kannan contraction need not be continuous. Indeed,
Kannan's condition controls the distance $d(Tx, Ty)$ only through the
displacements $d(x, Tx)$ and $d(y, Ty)$.
\end{remark}

\begin{theorem}[Kannan, 1968]\label{thm:kannan}
Let $(X, d)$ be a complete metric space and $T : X \to X$ a Kannan contraction
with constant $\alpha \in [0, 1/2)$. Then $T$ has a unique fixed point.
\end{theorem}

\begin{proof}
Let $x_0 \in X$ and $x_n = T^n(x_0)$. We have:
\[
    d(x_{n+1}, x_n) = d(Tx_n, Tx_{n-1})
    \leq \alpha \bigl[ d(x_n, Tx_n) + d(x_{n-1}, Tx_{n-1}) \bigr]
    = \alpha \bigl[ d(x_n, x_{n+1}) + d(x_{n-1}, x_n) \bigr].
\]
Hence $(1 - \alpha) \, d(x_{n+1}, x_n) \leq \alpha \, d(x_n, x_{n-1})$, so
\[
    d(x_{n+1}, x_n) \leq \frac{\alpha}{1 - \alpha} \, d(x_n, x_{n-1})
    = \beta \, d(x_n, x_{n-1})
\]
with $\beta = \alpha/(1 - \alpha) < 1$ since $\alpha < 1/2$.

Thus $d(x_{n+1}, x_n) \leq \beta^n d(x_1, x_0)$. The sequence $(x_n)$ is
Cauchy (same argument as for Banach) and converges to $x^* \in X$.

We verify that $x^*$ is a fixed point:
\[
    d(x^*, Tx^*) \leq d(x^*, x_{n+1}) + d(x_{n+1}, Tx^*)
    = d(x^*, x_{n+1}) + d(Tx_n, Tx^*)
    \leq d(x^*, x_{n+1}) + \alpha[d(x_n, x_{n+1}) + d(x^*, Tx^*)].
\]
Hence $(1 - \alpha) \, d(x^*, Tx^*) \leq d(x^*, x_{n+1}) + \alpha \, d(x_n, x_{n+1})
\to 0$, and $Tx^* = x^*$.

Uniqueness: if $x^*, y^*$ are two fixed points, then
$d(x^*, y^*) = d(Tx^*, Ty^*) \leq \alpha[d(x^*, Tx^*) + d(y^*, Ty^*)] = 0$.
\end{proof}

\begin{intuition}[title=\textbf{Kannan vs.\ Banach}]
Banach and Kannan contractions are independent classes:
\begin{itemize}
    \item There exist Banach contractions that are not Kannan contractions;
    \item There exist Kannan contractions that are not Banach contractions (and
          that are not even continuous).
\end{itemize}
Nevertheless, both conditions guarantee existence and uniqueness of a fixed
point in a complete metric space.
\end{intuition}

\begin{example}[Discontinuous Kannan map]
Let $X = [0, 1]$ and $T$ defined by $T(x) = x/4$ for $x \in [0, 1)$ and
$T(1) = 1/6$. Then $T$ is a Kannan contraction (with suitable $\alpha$) but $T$
is not continuous at $x = 1$. Its unique fixed point is $x^* = 0$.
\end{example}

\section{Chatterjea contractions}

\begin{definition}[Chatterjea mapping]
A mapping $T : X \to X$ is a \emph{Chatterjea contraction} if there exists
$\beta \in [0, 1/2)$ such that
\[
    d(Tx, Ty) \leq \beta \bigl[ d(x, Ty) + d(y, Tx) \bigr] \quad \forall x, y \in X.
\]
\end{definition}

\begin{theorem}[Chatterjea, 1972]
Let $(X, d)$ be a complete metric space and $T : X \to X$ a Chatterjea
contraction. Then $T$ has a unique fixed point.
\end{theorem}

\begin{proof}
Let $x_n = T^n(x_0)$. We have:
\begin{align*}
    d(x_{n+1}, x_n) &= d(Tx_n, Tx_{n-1})
    \leq \beta[d(x_n, Tx_{n-1}) + d(x_{n-1}, Tx_n)] \\
    &= \beta[d(x_n, x_n) + d(x_{n-1}, x_{n+1})]
    = \beta \, d(x_{n-1}, x_{n+1}) \\
    &\leq \beta[d(x_{n-1}, x_n) + d(x_n, x_{n+1})].
\end{align*}
Hence $(1 - \beta) d(x_{n+1}, x_n) \leq \beta \, d(x_n, x_{n-1})$, and we
conclude as for Kannan with the ratio $\beta/(1-\beta) < 1$.

The convergence of $(x_n)$ to a fixed point $x^*$ and uniqueness are proved
analogously.
\end{proof}

\section{\'Ciri\'c's contraction}

\begin{definition}[\'Ciri\'c quasi-contraction]
A mapping $T : X \to X$ is a \emph{\'Ciri\'c quasi-contraction} if there exists
$q \in [0, 1)$ such that
\[
    d(Tx, Ty) \leq q \cdot \max\{d(x,y),\, d(x, Tx),\, d(y, Ty),\, d(x, Ty),\, d(y, Tx)\}
\]
for all $x, y \in X$.
\end{definition}

\begin{remark}
\'Ciri\'c's condition simultaneously generalizes those of Banach ($q \cdot d(x,y)$),
Kannan ($q \cdot [d(x,Tx) + d(y,Ty)]/2$), and Chatterjea
($q \cdot [d(x,Ty) + d(y,Tx)]/2$).
\end{remark}

\begin{theorem}[\'Ciri\'c, 1974]\label{thm:ciric}
Let $(X, d)$ be a complete metric space and $T : X \to X$ a \'Ciri\'c
quasi-contraction with constant $q \in [0, 1)$. Then $T$ has a unique fixed
point $x^*$, and for every $x_0 \in X$, $T^n(x_0) \to x^*$.
\end{theorem}

\begin{proof}
Set $x_n = T^n(x_0)$ and $d_n = d(x_n, x_{n+1})$. We have:
\[
    d_n = d(Tx_{n-1}, Tx_n) \leq q \cdot \max\{d_{n-1}, d_{n-1}, d_n, d(x_{n-1}, x_{n+1}), 0\}.
\]
By the triangle inequality, $d(x_{n-1}, x_{n+1}) \leq d_{n-1} + d_n$, so
\[
    d_n \leq q \cdot \max\{d_{n-1}, d_n, d_{n-1} + d_n\}.
\]
If $d_n > d_{n-1}$, then $\max = d_{n-1} + d_n \leq 2d_n$, giving
$d_n \leq 2q d_n$, impossible for $q < 1/2$. A more delicate argument handles
the general case and shows $d_n \leq q \cdot d_{n-1}$ when $d_n \leq d_{n-1}$.

The sequence is Cauchy and converges to $x^*$. To verify $x^*$ is a fixed
point:
\[
    d(Tx^*, x^*) \leq d(Tx^*, Tx_n) + d(x_{n+1}, x^*)
    \leq q \cdot \max\{d(x^*, x_n), d(x^*, Tx^*), d_n, d(x^*, x_{n+1}), d(x_n, Tx^*)\} + d(x_{n+1}, x^*).
\]
Passing to the limit, $d(Tx^*, x^*) \leq q \cdot d(x^*, Tx^*)$, hence
$d(Tx^*, x^*) = 0$.

Uniqueness: if $y^* = Ty^*$, then $d(x^*, y^*) = d(Tx^*, Ty^*) \leq
q \cdot d(x^*, y^*)$, hence $d(x^*, y^*) = 0$.
\end{proof}

\section{Caristi's theorem}

Caristi's theorem is remarkable because it assumes no continuity condition on $T$.

\begin{theorem}[Caristi, 1976]\label{thm:caristi}
Let $(X, d)$ be a complete metric space and $\varphi : X \to [0, \infty)$ a
lower semicontinuous function. If $T : X \to X$ satisfies
\[
    d(x, Tx) \leq \varphi(x) - \varphi(Tx) \quad \forall x \in X,
\]
then $T$ has a fixed point.
\end{theorem}

\begin{proof}
Define a partial order on $X$ by
\[
    x \preceq y \iff d(x, y) \leq \varphi(x) - \varphi(y).
\]
We verify this is indeed a partial order:
\begin{itemize}
    \item Reflexivity: $d(x,x) = 0 \leq \varphi(x) - \varphi(x) = 0$.
    \item Antisymmetry: if $x \preceq y$ and $y \preceq x$, then $d(x,y) = 0$.
    \item Transitivity: if $x \preceq y$ and $y \preceq z$, then
          $d(x,z) \leq d(x,y) + d(y,z) \leq [\varphi(x) - \varphi(y)] + [\varphi(y) - \varphi(z)] = \varphi(x) - \varphi(z)$.
\end{itemize}

The hypothesis says $x \preceq Tx$ for every $x$.

By Zorn's lemma, it suffices to show every chain $(x_\alpha)_{\alpha \in A}$ in
$(X, \preceq)$ has an upper bound. Let $(x_\alpha)$ be a totally ordered chain.
The function $\alpha \mapsto \varphi(x_\alpha)$ is decreasing and bounded below
by $0$, so the net $(\varphi(x_\alpha))$ has an infimum $\ell$. One can extract
a generalized Cauchy subnet: for $\alpha \leq \beta$,
$d(x_\alpha, x_\beta) \leq \varphi(x_\alpha) - \varphi(x_\beta) \to 0$.
By completeness, it converges to $\bar{x}$. By lower semicontinuity,
$\varphi(\bar{x}) \leq \ell$. One verifies that $\bar{x}$ is an upper bound of
the chain.

By Zorn's lemma, there exists a maximal element $x^*$. Since $x^* \preceq Tx^*$
and $x^*$ is maximal, we have $Tx^* = x^*$.
\end{proof}

\begin{keyformulas}[title=\textbf{Caristi's condition}]
$d(x, Tx) \leq \varphi(x) - \varphi(Tx)$ with $\varphi$ l.s.c.\ and
$\varphi \geq 0$ in $(X, d)$ complete $\implies$ $T$ has a fixed point.
\end{keyformulas}

\section{Equivalence of Caristi and Ekeland's variational principle}

\begin{theorem}[Ekeland's variational principle, 1974]
Let $(X, d)$ be a complete metric space and $\varphi : X \to (-\infty, +\infty]$
lower semicontinuous, bounded below, and not identically $+\infty$. For every
$\varepsilon > 0$ and every $x_0 \in X$ with
$\varphi(x_0) \leq \inf_X \varphi + \varepsilon$, there exists $x^* \in X$
such that:
\begin{enumerate}
    \item $\varphi(x^*) \leq \varphi(x_0)$;
    \item $d(x_0, x^*) \leq 1$;
    \item $\varphi(x^*) < \varphi(x) + \varepsilon \, d(x, x^*)$ for all $x \neq x^*$.
\end{enumerate}
\end{theorem}

\begin{proposition}
Caristi's theorem and Ekeland's variational principle are equivalent.
\end{proposition}

\begin{proof}[Sketch]
\textbf{Ekeland $\Rightarrow$ Caristi}: If $T$ satisfies Caristi's condition,
let $x^*$ be given by Ekeland with suitable $\varepsilon$. Condition (3) with
$x = Tx^*$ and Caristi's condition lead to $Tx^* = x^*$.

\textbf{Caristi $\Rightarrow$ Ekeland}: More technical, using the construction
of a decreasing sequence for $\varphi$ in an appropriate ball.
\end{proof}

\section{Zamfirescu contractions}

\begin{definition}[Zamfirescu contraction]
$T : X \to X$ is a \emph{Zamfirescu contraction} if there exist
$a \in [0,1)$, $b \in [0, 1/2)$, $c \in [0, 1/2)$ such that for all $x, y$,
at least one of the following holds:
\begin{enumerate}
    \item $d(Tx, Ty) \leq a \, d(x, y)$;
    \item $d(Tx, Ty) \leq b[d(x, Tx) + d(y, Ty)]$;
    \item $d(Tx, Ty) \leq c[d(x, Ty) + d(y, Tx)]$.
\end{enumerate}
\end{definition}

\begin{theorem}[Zamfirescu, 1972]
Every Zamfirescu contraction on a complete metric space has a unique fixed point.
\end{theorem}

\begin{proof}
Set $\delta = \max\{a, b/(1-b), c/(1-c)\}$. Then $\delta < 1$ and
$d(Tx, Ty) \leq \delta \, d(x,y)$ for all $x, y$, reducing to Banach's theorem.

Indeed, if condition (1) holds, $d(Tx,Ty) \leq a \, d(x,y)$. If condition (2)
holds, $d(Tx,Ty) \leq b[d(x,Tx) + d(y,Ty)] \leq b[d(x,y) + 2d(y, Ty)] + b \, d(Tx, Ty)$,
giving $(1-b) d(Tx,Ty) \leq b[d(x,y) + 2d(y,Ty)]$, and careful reasoning shows
$d(Tx,Ty) \leq \frac{2b}{1-b} d(x,y)$. A similar argument applies for (3).
\end{proof}

\section{Subrahmanyam's characterization}

\begin{theorem}[Subrahmanyam, 1975]
A metric space $(X, d)$ is complete if and only if every Kannan contraction
$T : X \to X$ has a fixed point.
\end{theorem}

\begin{remark}
This result is remarkable because it shows that Kannan's condition
\emph{characterizes} metric completeness, unlike Banach's condition. There exist
incomplete metric spaces in which every Banach contraction has a fixed point.
\end{remark}

\section{Exercises}

\begin{exercise}
Show that $T(x) = \sin(x)$ is a weak contraction on $[0, \pi]$ but not a
Banach contraction. Apply Edelstein's theorem.
\end{exercise}

\begin{exercise}
Construct a Kannan contraction that is not continuous and verify directly that
it has a unique fixed point.
\end{exercise}

\begin{exercise}
Show that every Banach contraction is a \'Ciri\'c quasi-contraction, but the
converse is false.
\end{exercise}

\begin{exercise}
Apply Caristi's theorem to the operator $T : C([0,1]) \to C([0,1])$ defined by
$(Tf)(t) = \frac{1}{2} f(t)$, with $\varphi(f) = \norm{f}_\infty$.
\end{exercise}

\begin{exercise}
Prove in detail the equivalence between Caristi's theorem and Ekeland's
variational principle.
\end{exercise}

\begin{exercise}
Let $(X, d)$ be complete and $T : X \to X$ such that for every $\varepsilon > 0$,
there exists $\delta > 0$ with $\varepsilon \leq d(x,y) < \varepsilon + \delta
\Rightarrow d(Tx, Ty) < \varepsilon$ (Meir--Keeler condition). Show that this
condition is stronger than Edelstein's weak contraction but weaker than Banach
contraction.
\end{exercise}
